\clearpage
\VerdierPhysicalPageAnchor{0060}%
\VerdierSemanticAnchor{ch1-1-5-14-continuation}
然后令：

\VerdierSemanticAnchor{ch1-1-5-14-eq3}
\begin{equation}
  \begin{aligned}
    T(0)&=\operatorname{id}_{\mathcal C},\\
    T(n)&=T^n, && n>0,\\
    T(n)&=(T^{-1})^{-n}, && n<0.
  \end{aligned}
  \tag{1.5.14.3}
\end{equation}

\VerdierSemanticAnchor{ch1-1-5-14-cocycle-data}
再对每一对整数 $(n,n')$，如下定义同构
\[
  c(n',n)\colon T(n+n')\longrightarrow T(n)T(n')
\]
：

\VerdierSemanticAnchor{ch1-1-5-14-eq4}
\begin{equation}
  \begin{aligned}
    n,n'&\geq 0,
      &c(n',n)&=\operatorname{id},\\
    n,n'&\leq 0,
      &c(n',n)&=\operatorname{id},\\
    n'&\leq 0, n\geq 0, n+n'\leq 0,
      &c(n',n)&=u^{(n)}\mathbin{*}T(n+n'),\\
    n'&\leq 0, n\geq 0, n+n'\geq 0,
      &c(n',n)&=T(n+n')\mathbin{*}u^{(-n')},\\
    n'&\geq 0, n\leq 0, n+n'\leq 0,
      &c(n',n)&=T(n+n')\mathbin{*}v^{(n')},\\
    n'&\geq 0, n\leq 0, n+n'\geq 0,
      &c(n',n)&=v^{(-n)}\mathbin{*}T(n+n').
  \end{aligned}
  \tag{1.5.14.4}
\end{equation}

\VerdierSemanticAnchor{ch1-1-5-14-conclusion}
为了确认这样确实定义了一个 $\mathbb Z$-范畴，还需验证 (1.2.1.3) 型的
恒等式成立。该验证建立在态射 $u^{(n)}$ 与 $v^{(n)}$ 所满足的经典恒等式上。
我们把细节留给读者。

\VerdierSemanticAnchor{ch1-1-6}
\subsection{Koszul 余循环与符号规则}

\VerdierSemanticAnchor{ch1-1-6-1}
\noindent\textbf{1.6.1.}\quad
设 $e_1,\ldots,e_n$ 为 $\mathbb Z^n$ 的一组基。对
$i,j\in[1,n]$ 且 $i\neq j$，以 $\mathbb Z^n_{i,j}$ 表示由 $e_i$ 与
$e_j$ 生成的 $\mathbb Z^n$ 的子群。限制同态给出同态

\VerdierSemanticAnchor{ch1-1-6-1-eq1}
\begin{equation}
  H^2(\mathbb Z^n,\mathbb Z/2\mathbb Z)
  \longrightarrow
  \bigoplus_{\substack{i,j\in[1,n]\\ i<j}}
    H^2(\mathbb Z^n_{i,j},\mathbb Z/2\mathbb Z).
  \tag{1.6.1.1}
\end{equation}

我们有（典范）同构

\VerdierSemanticAnchor{ch1-1-6-1-eq2}
\begin{equation}
  H^2(\mathbb Z^n_{i,j},\mathbb Z/2\mathbb Z)
  \simeq \mathbb Z/2\mathbb Z.
  \tag{1.6.1.2}
\end{equation}

\VerdierSemanticAnchor{ch1-1-6-1-conclusion}
同态 (1.6.1.1) 是同构。
